\documentclass{article}
\usepackage{enumerate}
\usepackage{amssymb}
\usepackage{amsmath}
\usepackage{amsthm}
\usepackage{latexsym}
\usepackage[T1]{fontenc}
\usepackage[latin1]{inputenc}


% JDLM headers

\usepackage{fancyhdr}
\pagestyle{fancy}

\let\titlepageAuthor\author
\renewcommand{\author}[1]{\newcommand{\headerAuthor}{#1}\titlepageAuthor{#1}} 
\let\titlepageTitle\title
\renewcommand{\title}[1]{\newcommand{\headerTitle}{#1}\titlepageTitle{#1}} 

\lhead{\headerTitle: \leftmark} \chead{} \rhead{\thepage} 
\lfoot{\copyright \headerAuthor} \cfoot{} \rfoot{ - MathNerds Course Guide Collection - www.jiblm.org}
\renewcommand{\headrulewidth}{0.4pt}
\renewcommand{\footrulewidth}{0.4pt}


\begin{document}


\title{Jordan Curve Theorem}
\author{William S. Mahavier}
\date{ }
\maketitle
\thispagestyle{empty}



\newpage

\pagenumbering{arabic}

\section*{Introduction}
\markboth{Introduction}{Introduction}

This section was developed by Wm. S. Mahavier, who adapted the material from the text, (forthcoming). It is a good example of how one can take a traditional book and break the material up into bite-sized pieces for students to absorb. The sequence is probably best for students who have had at least an introduction to topology, and is certainly appropriate at the graduate level.

\section*{Theorem Sequence}
\markboth{Theorem Sequence}{Theorem Sequence}

All point sets in this section are subsets of $\Re^2$. A rectangular grating is the union of a square with sides parallel to the axes, and a finite collection of line segments each of which is either vertical or horizontal and has both end points on the square. The $2$-cells of a rectangular grating ${\cal G}$ are the closures of the components of $\Re^2 - {\cal G}$ (one of which is not really a $2$-cell.) The $1$-cells of ${\cal G}$ are the sides of the bounded $2$-cells of ${\cal G}$. The $0$-clells of ${\cal G}$ are the corners of the bounded $2$-cells of ${\cal G}$.

\begin{description}

  \item[Theorem 1:] If $H$ and $K$ are disjoint closed subsets of $\Re^2$ and $H$ is bounded, then there is a grating ${\cal G}$ such that no $2$-cell of ${\cal G}$ intersects both $H$ and $K$.

  \item[] A $k$-chain on a grating ${\cal G}$ is a function from the set of $k$-cells of ${\cal G}$ into the set $\{ 0,1 \}$. Obviously this is equivalent to choosing a subcollection of the $k$-cells of ${\cal G}$. If each of $C$ and $D$ is a $k$-chain, then $C+D$ is the $k$-chain such that $(C+D)(M) = 0$ if and only if $C(M) = D(M) = 1$ or $C(M) = D(M) = 0$.

  \item[] The $k$-chain of a grating ${\cal G}$ with this operation are denoted $C_k ({\cal G})$. The $k$-chain which is $1$ only at the $k$-cell $K$ will be denoted $\tilde{K}$.

  \item[Theorem 2:] The $k$-chains on a grating ${\cal G}$ form a commutative group.

  \item[Theorem 3:] There is a homomorphism $\delta_2$ from $C_2({\cal G})$ into $C_1({\cal G})$ such that if $K$ is a $2$-cell of the grating ${\cal G}$, then $\delta_2(\tilde{K})$ is $1$ only at the $1$-cells which are subsets of $K$. Moreover, if $C$ is a $2$-chain and $L$ is a $1$-cell, then $\delta_2(C)(L) = 1$ if and only if there is an odd number of $2$-cells $K$ such that $C(K)=1$ and $K$ contains $L$.

  \item[Theorem 4:] There is a homomorphism $\delta_1$ from $C_1({\cal G})$ into $C_0({\cal G})$ such that if $K$ is a $1$-cell of the grating ${\cal G}$, then $\delta_1(\tilde(K)$ is $1$ only at the $0$-cells which are subsets of $K$. Moreover, if $C$ is a $1$-chain and $L$ is a $0$-cell, then $\delta_1(C)(L)=1$ if and only if there is an odd number of $1$-cells $K$ such that $C(K)=1$ and $K$ contains $L$.

  \item[] We will ignore the subscript and denote both $\delta_2$ and $\delta_1$ by $\delta$. These homomorphisms are called boundary operators, and for the $k$-chain $C$, $\delta C$ is called the boundary of $C$.

  \item[] The statement that a $k$-chain $C$ is a cycle means that $k=0$ or $k=1$ and $\delta C = 0$ or $k=2$ and $\delta C = 0$.

  \item[Theorem 5:] The $k$-cycles form a subgroup of $C_k({\cal G})$.

  \item[Theorem 6:] If $k$ is positive and $C$ is a $k$-chain, then $\delta C$ is a $(k-1)$-cycle.

  \item[] A cycle is said to be a bounding cycle if it is in the boundary of a chain.

  \item[Theorem 7:] If $C$ is a bounding $0$-cycle, there are an even number of $0$-cells $K$ such that $C(K) = 1$.

  \item[Theorem 8:] If $D$ is a $1$-cycle and $|D|$ is finite, then $\ldots \ni$ ???

  \item[] If $C$ is a $k$-chain, the carrier of $C$, denoted $|C|$, is the union of all $k$-cells $K$ such that $C(K)=1$. A $k$-chain $C$ is said to be connected if $|C|$ is connected. A $k$-chain $D$ is said to be a component of a $k$-chain $C$ if $D$ is connected and $|D|$ is a component of $|C|$. A $k$-chain $C$ is said to be in a point set $M$ is $|C|$ is a subset of $M$.

  \item[Theorem 9:] If $D$ is a component of a $k$-chain $C$ and $k$ is positive, then $|\delta D|$ is $|\delta C| \cap |D|$.

  \item[Theorem 10:] If $C$ is a $1$-chain and $|\delta C|$ is a set consisting of two points $p$ and $q$, then $p$ and $q$ are in the same component of $|C|$.

  \item[Theorem 11:] If $C$ is a $2$-chain, $|\delta C|$ is the point set boundary of $|C|$.

  \item[] If each of ${\cal G}$ and $H$ is a grating, then $H$ is said to be a refinement of ${\cal G}$ if $H$ contains ${\cal G}$.

  \item[Theorem 12:] If ${\cal G}$ and $H$ are gratings, $H$ is a refinement of ${\cal G}$, and $C$ is a $k$-chain on ${\cal G}$, then there is only one $k$-chain $D$ on $H$ such that $|C| = |D|$.

  \item[Definition 13:] Suppose each of ${\cal G}$ and $H$ is a grating, and $H$ is a refinement of ${\cal G}$. For each $k$-chain $C$ on ${\cal G}$, $sdC$ denotes the $c$-chain on $H$ such that $|sdC| = |C|$. The chain $sdC$ is called a \emph{subdivision} of $C$.

  \item[Theorem 14:] Each $1$-cycle is the boundary of exactly two $2$-chains.

  \item[] Suppose $M$ is a point set and $C$ is a $k$-cycle on a grating ${\cal G}$. Then $C$ bounds in $M$ means that there is a refinment $H$ of ${\cal G}$, and a $(k+1)$-chain $D$ on $H$ such that $|D|$ is a subset of $M$ and $\delta D = sdC$.

  \item[Theorem 15:] Suppose each of $C_1, C_2, \cdots, C_2$ is a $k$-cycle on the grating ${\cal G}$ which bounds in the points set $M$. Then $C_1 + C_2 + \cdots + C_n$ bounds in $M$.

  \item[Theorem 16:] If the $k$-cycle $C$ does not bound in the points set $M$, then some component of $C$ does not bound in $M$.

  \item[Theorem 17:] If $\Re^2-M$ is connected, then every $1$-cycle in $M$ bounds in $M$.

  \item[Lemma 18:] Suppose $U$ is an open set and $p$ is a point of $U$. The set of all points $q$ of $U$ such that there is a grating ${\cal G}$ and a $1$-chain on ${\cal G}$ in $U$ whose boundary is $\tilde{p} + \tilde{q}$, is both open and closed in $U$.

  \item[Theorem 19:] Suppose $U$ is an open set, $p$ and $q$ are two points of $U$, $p$ and $q$ are $0$-cells of the grating ${\cal G}$, and $C$ is $\tilde{p} + \tilde{q}$. Then $C$ bounds in $U$ if and only if $p$ and $q$ are in the same component of $U$.

  \item[Theorem 20:] If $C$ and $D$ are $1$-cycles and $p$ and $q$ are two points not in $|C| \cup |D|$, then at least one of the cycles $C$, $D$, and $C+D$ bounds in $\Re^2 - \{ p,q \}$.

  \item[] Two $1$-chains $C$ and $D$ on a grating ${\cal G}$ are said to have general intersection if
    \begin{enumerate}[\hspace{1cm}(1)]
      \item no $0$-cell common to $|C|$ and $|D|$ is the outer edge of ${\cal G}$, and
      \item at each $0$-cell $p$ common to $|C|$ and $|D|$, the horizontal $1$-cells of ${\cal G}$ containing $p$ are contained in only one of $|C|$ and $|D|$ and the vertical $1$-cells of ${\cal G}$ containing $p$ are contained in only one of $|C|$ and $|D|$.
    \end{enumerate}

  \item[Theorem 21:] Suppose $C$ is a $1$-cycle and $D$ is a $1$-chain having general intersection with $C$ and whose boundary is $\tilde{p} + \tilde{q}$. Then $C$ bounds in $\Re^2 - \{p,q\}$ if and only if $|C|$ and $|D|$ intesect at an even number of $0$-cells.

  \item[Theorem 22:] If two $1$-cycles $C$ and $D$ have general intersection, then $|C|$ and $|D|$ intersect at an even number of $0$-cells.

  \item[Theorem 23:] Suppose $C$ and $D$ are $1$-chains having general intersection, $\delta C = \tilde{p} + \tilde{q}$, and $\delta D = \tilde{r} + \tilde{s}$, $|C| \cap |D|$ contains an odd number of $0$-cells and $M$ is a continuum which contains $p$ and $q$ but does not intersect $|D|$. Then $M \cup |C|$ separates $r$ and $s$.

  \item[Lemma 24:] Suppose $U$ and $V$ are open sets, $U$ is bounded, ${\cal G}$ is a grating, and $C$ is a chain on ${\cal G}$ such that $|C| \subseteq U \cup V$. Then there is a refinement $H$ of ${\cal G}$ such that every cell of $H$ which is contained in $|sdC|$ is either a subset of $U$ or a subset of $V$.

  \item[Theorem 25:] Suppose $U$ and $V$ are open sets in $\Re^2$; $U$ is bounded; $C$ and $D$ are $1$-chains such that $C$ is in $U$, $D$ is in $V$, and $\delta C = \delta D = \tilde{p} + \tilde{q}$; and $C+D$ bounds in $U \cup V$. Then $\tilde{p} + \tilde{q}$ bounds in $U \cap V$.

  \item[Theorem 26:] Suppose $M$ and $N$ are closed subsets of $\Re^2$ and $M$ is bounded or $M$ and $N$ are disjoint. If $\delta C = \delta D = \tilde{p} + \tilde{q}$, $|C|$ does not intersect $M$, and $|D|$ does not intersect $N$ but $C+D$ bounds in $\Re^2 - (M \cap N)$, then $p$ and $q$ are not separated by $M \cup N$.

  \item[Theorem 27:] Suppose that $M$ and $N$ are closed subsets of $\Re^2$, and either $M$ and $N$ are disjoint or $M$ is bounded and $M \cap N$ is connected. If $p$ and $q$ are not separated by either $M$ or $N$, then $P$ and $q$ are not separated by $M \cup N$.

  \item[Theorem 28:] If $U$ and $V$ are connected open sets whose union is $\Re^2$, then the intersection of $U$ and $V$ is connected.

  \item[Theorem 29:] Suppose $M$ is a closed set contained in a connected open set $U$ in $\Re^2$ and $V_1, V_2, \cdots$ are components of the point set $\Re^2-M$. Then the components of $U-M$ are $U \cap V_1, U \cap V_2, \cdots$.

  \item[Theorem 30:] Suppose $U$ and $V$ are open subsets of $\Re^2$ which do not separate $\Re^2$, $K$ and $L$ are $1$-chains such that $|K|$ is a subset of $U$ and $|L|$ is a subset of $V$, and $\delta K = \delta L$. If $\delta K$ bounds in $U \cap V$, then $K+L$ bounds in $U \cup V$.

  \item[Theorem 31:] Suppose $M$ and $N$ are connected closed subset of $\Re^2$ and $C$ is a $0$-cycle which bounds a $1$-chain $K$ such that $|K|$ does not intersect $M$ and bounds a $1$-chain $L$ such that $|L|$ does not intersect $N$. Suppose the $2n$ $0$-cells of $|C|$ are listed in some way as $(x_1,y_1), (x_2,y_2), \cdots, (x_n,y_n)$. If $K+L$ does not bound in $\Re^2-(M \cap N)$, then for some $r$, $M \cup N$ separates $x_r$ from $y_r$.

  \item[Theorem 32:] No arc separates $\Re^2$.

  \item[Theorem 33:] If $J$ is a simple closed curve in $\Re^2$, then $\Re^2-J$ is the union of two open sets each of which has $J$ as its point set boundary.

  \item[Theorem 34:] If each of $M$ and $N$ is a closed connected subset of $\Re^2$ and $M$ is bounded, but $M \cap N$ is not connected, there is a pair of points separated by $M \cup N$ but not by $M$.

  \item[Theorem 35:] The unit $2$-cell in $\Re^2$ is unicoherent; that is, if $M$ and $N$ are closed connected sets whose union is $\Re^2$, then $M \cap N$ is connected.

\end{description}

\end{document}
